\subsection{An example from complex Gram coordinates}
\label{sec:gram-composition-constructions}

The following sextic is obtained from the nine quadratic coordinates
of a complex Hermitian Gram matrix. Its profile uses the imaginary
part of a cyclic product of three off-diagonal entries. It shows
that the composition criterion also produces higher-degree minimal
regular zero sets.

\begin{example}[A complex Gram sextic]
\label{thm:complex-gram-cycle-sextic}
For every $k\ge2$, let $x_1,x_2,x_3\in\mathbb C^k$ and
$h_{ij}=x_i^*x_j$.  Then
\begin{equation}\label{eq:complex-gram-cycle-sextic}
 F=\operatorname{Im}(h_{12}h_{23}h_{31})
\end{equation}
is a real harmonic homogeneous polynomial of degree six on $\R^{6k}$
with nonempty minimal regular zero set.

More explicitly, set $r_i=h_{ii}$, $R=\sum_i r_i$,
\[
 a_1=|h_{23}|^2,\qquad a_2=|h_{31}|^2,\qquad
 a_3=|h_{12}|^2,\qquad
 t=\operatorname{Re}(h_{12}h_{23}h_{31}).
\]
With $\{i,j,l\}=\{1,2,3\}$ in each summand, define
\begin{align}
 S&=\sum_{i=1}^3 r_i a_i(a_j+a_l)-2t\sum_{i=1}^3a_i,
 \label{eq:complex-gram-cycle-gradient}\\
 B&=\sum_{i=1}^3a_i(r_i+r_j)(r_i+r_l)
       -\sum_{i=1}^3a_i^2-2Rt.
 \label{eq:complex-gram-cycle-weight}
\end{align}
Then
\begin{equation}\label{eq:complex-gram-cycle-minimality}
 \Delta F=0,\qquad |\nabla F|^2=S,\qquad
 \Gamma(F,S)=4BF,\qquad \cL(F)=-2BF.
\end{equation}
\end{example}

\begin{proof}
Use the complex bilinear extension of $\Gamma$ in
\eqref{eq:complex-Gram-differential-table}.  For
$Z=h_{12}h_{23}h_{31}$ and $A=a_1+a_2+a_3$, the product rule gives
\[
 \Gamma(Z,Z)=4AZ,\qquad
 \Gamma(Z,\overline Z)=2\sum_i r_i a_i(a_j+a_l).
\]
Taking real and imaginary parts proves
$|\nabla F|^2=S$.  The same table gives
\begin{equation}\label{eq:complex-gram-cycle-pairings}
 \Gamma(F,r_i)=4F,\quad
 \Gamma(F,a_i)=2(R-r_i)F,\quad
 \Gamma(F,t)=2AF.
\end{equation}
Applying these three identities to the expression for $S$ gives
$\Gamma(F,S)=4BF$. Indeed, writing
$D=\sum_i r_i a_i$ and $E=\sum_i r_i a_i^2$, one has
$S=AD-E-2tA$. After the $A^2$ terms cancel, the product rule gives
\[
 \Gamma(F,S)=4F\left[
 RD+\frac A2\sum_i r_i(R-r_i)-\sum_i a_i^2
 -\sum_i r_i a_i(R-r_i)-2Rt\right].
\]
The coefficient of $a_i$ in the brackets is
$r_i^2+r_1r_2+r_2r_3+r_3r_1=(r_i+r_j)(r_i+r_l)$,
so the brackets equal $B$. This is a polynomial identity, valid
also when the Gram entries satisfy rank relations, as for $k=2$.
For harmonicity, with $x_2,x_3$ fixed, $F$ is the Hermitian
quadratic form in $x_1$ with matrix
\[
 \frac{h_{23}x_2x_3^*-\overline{h_{23}}x_3x_2^*}{2i}.
\]
Its trace is zero, so its block Laplacian vanishes.  The same
argument applies to the other two blocks.  The point
$x_1=x_3=e_1$, $x_2=e_1+e_2$ has $F=0$ and $S=2$, proving
nonempty regularity; \eqref{eq:complex-gram-cycle-minimality}
then proves minimality.

\end{proof}

